\documentclass[12pt,a4paper]{article}

\usepackage{fontspec}
\setmainfont{TeX Gyre Pagella}
\usepackage{amsmath,amssymb,amsthm}
\usepackage[margin=1.1in]{geometry}
\usepackage{microtype}
\usepackage[colorlinks=true,linkcolor=black,citecolor=black,urlcolor=black]{hyperref}
\usepackage{algorithm}
\usepackage{algpseudocode}
\usepackage{booktabs}
\usepackage[shortlabels]{enumitem}
\usepackage{setspace}
\usepackage{titlesec}
\usepackage{parskip}
\setlength{\parindent}{0pt}
\setlength{\parskip}{0.8em}

\onehalfspacing

\theoremstyle{plain}
\newtheorem{theorem}{Theorem}[section]
\newtheorem{proposition}[theorem]{Proposition}
\newtheorem{lemma}[theorem]{Lemma}
\newtheorem{corollary}[theorem]{Corollary}

\theoremstyle{definition}
\newtheorem{definition}[theorem]{Definition}
\newtheorem{example}[theorem]{Example}
\newtheorem{notation}[theorem]{Notation}

\theoremstyle{remark}
\newtheorem{remark}[theorem]{Remark}
\newtheorem{hypothesis}[theorem]{Empirical Hypothesis}
\newtheorem{proposal}[theorem]{Architectural Proposal}

\DeclareMathOperator{\Disp}{Disp}
\DeclareMathOperator{\Obs}{Obs}
\DeclareMathOperator{\Adm}{Adm}

\title{\textbf{Cinematic Reconstruction}\\[4pt]
\large A Framework for Video Generation, Compression, and Diagnosis\\
via Persistent World State}
\author{Flyxion\\[2pt]
\small Independent Researcher}
\date{2026}

\begin{document}

\maketitle
\thispagestyle{empty}

\begin{abstract}
We develop a formal framework treating video not as a sequence of independently synthesized images but as a sequence of camera observations of a persistent, evolving world state. A movie is then fundamentally a family of partial observations of a latent system whose full structure exceeds what any single frame can reveal. This perspective reorganizes several problems that resist solution in frame-local approaches: occlusion and reappearance, temporal consistency, style transfer as realization rather than regeneration, camera movement as observation-operator change, generative quality control, and compression artifact localization. The framework incorporates a heterogeneous algorithm repertoire in which multiple independent procedures propose and evaluate candidate continuations, and a formal witness and disagreement mechanism that converts model disagreement from an inconvenience into diagnostic information. Glitches are defined as violations of preserved or recoverable structural invariants rather than as merely unattractive pixel configurations, enabling them to serve as evidence about which representation or algorithm failed and where. A structure-aware rate-distortion objective distinguishes degradation by type rather than treating all distortion as equivalent. An appendix treats the residual underdetermination of the rendering process as a selectable information channel, subordinate to admissibility, permitting embedding of signals in otherwise unconstrained scene and material choices. The framework draws on established results in multi-view geometry, SLAM, Gaussian splatting, neural radiance fields, video compression, rate-distortion theory, partial observability, ensemble uncertainty, and steganography, while proposing a particular architectural organization that places persistent world state at the center of all stages from generation through compression through diagnosis.
\end{abstract}

\newpage
\tableofcontents
\newpage

%% ============================================================
\section{Introduction}
%% ============================================================

Contemporary video generation systems and video codecs share a common implicit assumption: the fundamental unit of video is a frame. A generative model learns to synthesize plausible images and, by various forms of temporal conditioning, to synthesize plausible sequences. A video codec decomposes each frame into blocks, transforms and quantizes those blocks, and optionally exploits inter-frame redundancy through motion compensation and prediction. In both cases the frame remains primary and temporal structure is secondary.

This monograph argues that the assumption is architecturally limiting in a precise sense. A frame is a projection of a scene. It discards information that was available in the scene in order to form a two-dimensional image: depth, material identity behind occluded surfaces, the continuation of objects outside the frame boundary, the state of characters whose faces are not currently visible, the trajectory of a moving object at the moment it passes behind a tree. This information does not disappear from the world when it disappears from the frame. It should not disappear from a system attempting to model, generate, reconstruct, or compress the video.

The central proposal is: treat a movie as a sequence of camera observations of a persistent, evolving world state rather than as a sequence of independently synthesized or compressed images. This distinction carries formal consequences for every stage of video production and consumption, which the sections below develop.

The framework is not entirely novel in its components. Persistent world representations appear in SLAM \cite{davison2007monoslam}, in 3D scene reconstruction, in game engines, in simulation environments, and in recent neural radiance and Gaussian splatting approaches to novel-view synthesis \cite{mildenhall2020nerf,kerbl2023gaussians}. Ensemble disagreement as an uncertainty signal appears throughout machine learning. Multi-component rate-distortion objectives have been explored in perceptual image and video compression \cite{wang2004ssim}. Steganographic embedding in images is a mature field \cite{fridrich2009steganography}. The contribution of the present framework is not any single technique but their organization under a common admissible reconstruction architecture, in which persistent world state is primary, rendered video is a partial observation of that state, heterogeneous algorithms propose and witness candidate continuations, structural disagreement drives diagnosis, and residual admissible freedom in the rendering process can optionally be exploited as an information channel.

Section~\ref{sec:worldstate} formalizes the world state and observation map. Section~\ref{sec:admissible} develops admissible reconstruction and the probate mechanism. Section~\ref{sec:repertoire} describes the heterogeneous algorithm repertoire. Section~\ref{sec:witnesses} develops the witness and disagreement formalism. Section~\ref{sec:glitches} defines structural glitches and their diagnostic role. Section~\ref{sec:ratedistortion} compares conventional and structure-aware rate-distortion objectives. Sections~\ref{sec:style} through \ref{sec:camera} address style, narrative, and cinematography. Sections~\ref{sec:scene} and \ref{sec:provenance} specify scene architecture and provenance. Section~\ref{sec:architecture} provides the modular system specification. Section~\ref{sec:experiments} proposes concrete experiments. Section~\ref{sec:implementation} outlines a staged implementation plan. Section~\ref{sec:limitations} discusses limitations, and Section~\ref{sec:conclusion} concludes. Appendix~\ref{app:steg} treats hyperpleonastic steganography.

%% ============================================================
\section{World State and Frame Observation}
\label{sec:worldstate}
%% ============================================================

\subsection{The World State Tuple}

\begin{definition}[World State]
A \textbf{world state} at time $t$ is a tuple
\[
\mathcal{W}_t = (G_t,\; O_t,\; M_t,\; L_t,\; C_t,\; S_t,\; H_t)
\]
where the components are:
\begin{itemize}[leftmargin=2em,itemsep=0pt]
\item $G_t$: geometry and spatial structure (meshes, volumes, implicit fields, skeletal rigs, and any hybrid representations);
\item $O_t$: persistent object identities, their bounding regions, semantic types, and inter-object relationships;
\item $M_t$: material and appearance properties for each geometric region, including reflectance, subsurface parameters, and procedural specifications;
\item $L_t$: illumination, including light source positions, intensities, spectra, and environment maps;
\item $C_t$: camera and observation parameters, including intrinsics, extrinsics, lens models, and motion blur exposure intervals;
\item $S_t$: semantic, narrative, and causal state, covering character knowledge, relationships, event history, physical state, and story-level facts;
\item $H_t$: the provenance history of the current state, recording which processes proposed each committed update and what evidence supported the commitment.
\end{itemize}
\end{definition}

This definition is intentionally broad. The components $G_t$, $M_t$, and $L_t$ correspond roughly to the scene description in physically based rendering \cite{pharr2023pbr}. The component $O_t$ adds object identity in a representation-independent sense: an object remains the same object whether it is currently represented as a mesh, a Gaussian splat, a neural radiance field, or a voxel grid. The components $S_t$ and $H_t$ extend the scope beyond geometry to include semantic and causal structure and the evidentiary record of how the world reached its current state.

\begin{remark}
Not all components need be present simultaneously in every application. A system focused on geometric consistency may omit $S_t$; a system focused on narrative generation may treat $G_t$ as a compressed sketch. The tuple defines the full scope; implementations instantiate what is needed.
\end{remark}

\subsection{The Observation Map}

\begin{definition}[Observation Map]
Given a world state $\mathcal{W}_t$, camera parameters $C_t$, and rendering parameters $\theta_t$, the \textbf{rendered frame} is
\[
I_t = R(\mathcal{W}_t,\; C_t,\; \theta_t).
\]
When the camera is treated as part of the world state, we write simply
\[
I_t = \Obs_{C_t}(\mathcal{W}_t).
\]
A video sequence of length $n$ is then a sequence of observations
\[
(I_1, I_2, \ldots, I_n) = (\Obs_{C_1}(\mathcal{W}_1),\; \Obs_{C_2}(\mathcal{W}_2),\; \ldots,\; \Obs_{C_n}(\mathcal{W}_n)).
\]
\end{definition}

The rendering function $R$ may be a physical light transport simulator, a real-time rasterizer, a neural renderer, a style-conditioned diffusion model, or any composition of such procedures, provided it maps world state and camera to a frame. Its precise form is not fixed by the framework but is instead a parameter of each application.

\subsection{The Frame as Projection}

\begin{proposition}[Non-injectivity of Observation]
For any fixed camera configuration $C$ and rendering parameters $\theta$, the observation map $\Obs_C$ is generically non-injective: there exist distinct world states $\mathcal{W} \neq \mathcal{W}'$ such that $\Obs_C(\mathcal{W}) = \Obs_C(\mathcal{W}')$.
\end{proposition}

\begin{proof}[Justification]
This follows from simple dimensionality arguments. A rendered frame $I \in \mathbb{R}^{H \times W \times 3}$ has at most $3HW$ real degrees of freedom. The world state $\mathcal{W}$, containing geometry, object identities, materials, illumination, and semantic content, has far more degrees of freedom than any finite-resolution image can represent. Distinct world states related by occlusion, by out-of-frame content, or by perceptually equivalent material configurations will render to the same image.
\end{proof}

\begin{definition}[Observational Equivalence Class]
The \textbf{observational equivalence class} of a world state $\mathcal{W}$ under camera sequence $(C_1, \ldots, C_n)$ is
\[
[\mathcal{W}]_{C_{1:n}} = \left\{ \mathcal{W}' \;:\; \Obs_{C_t}(\mathcal{W}') = \Obs_{C_t}(\mathcal{W}) \text{ for all } t \in \{1, \ldots, n\} \right\}.
\]
\end{definition}

The observational equivalence class is generally non-trivial and shrinks as more camera observations are accumulated. This is the core information-theoretic fact underlying all reconstruction problems: observations constrain but generally do not determine the latent world.

\subsection{Occlusion as Observation Loss}

\begin{proposition}[Occlusion Does Not Delete]
If an object $o \in O_t$ is occluded by another object at time $t$ under camera $C_t$, the observation $I_t = \Obs_{C_t}(\mathcal{W}_t)$ contains no pixels from $o$, but $o$ remains a component of $O_t$ and of $\mathcal{W}_t$ with all its geometric, material, positional, and identity information intact.
\end{proposition}

This proposition is trivial given the definitions but has substantial architectural consequences. A system that maintains $\mathcal{W}_t$ explicitly can continue tracking object $o$ through occlusion, maintain its velocity and trajectory, and produce a geometrically consistent rendering when the camera moves to reveal the object again. A frame-local system that discards $\mathcal{W}_t$ and retains only $I_t$ loses this information.

\begin{hypothesis}
\label{hyp:occlusion}
A system maintaining a persistent world state $\mathcal{W}_t$ through occlusion events produces geometrically and photometrically more consistent object reappearance than a frame-local system conditioned only on recent frames, when measured on controlled scenes with known ground truth.
\end{hypothesis}

This is an empirical hypothesis to be tested; its experimental design is specified in Section~\ref{sec:experiments}.

\subsection{Camera Motion as Observation Operator Change}

Under the framework, camera motion changes the observation operator rather than recreating or regenerating the world. A pan, zoom, cut, rack focus, or crane shot corresponds to a change in $C_t$ within the world trajectory $(\mathcal{W}_t)_{t \geq 0}$. The world itself may or may not change simultaneously. When only the camera changes, the world state component $G_t, O_t, M_t, L_t$ may remain constant while the rendered image changes substantially. This is both the ordinary situation in cinematography and a point of significant difficulty for frame-local generative models, which must implicitly infer object permanence from a limited temporal context window rather than from an explicit world model.

%% ============================================================
\section{Admissible Reconstruction}
\label{sec:admissible}
%% ============================================================

\subsection{Reconstruction as Constrained Inference}

The inverse problem of video is: given an observation sequence $(I_1, \ldots, I_t)$, what can be recovered about the latent world trajectory $(\mathcal{W}_s)_{s \leq t}$? In general, this is ill-posed; the previous section established that the observation map is non-injective. The appropriate response is not to select one world hypothesis arbitrarily but to characterize the set of hypotheses consistent with the observations.

\begin{definition}[Admissible Reconstruction Set]
Given an observation sequence $(I_1, \ldots, I_t)$, a constraint system $\mathcal{C}$ (encoding physical laws, object permanence, causal constraints, semantic consistency, and user-specified facts), and a history $H_t$, the \textbf{admissible reconstruction set} is
\[
\mathcal{R}(I_{1:t},\; \mathcal{C},\; H_t) = \left\{ \mathcal{W} \;:\; \Obs_{C_s}(\mathcal{W}) = I_s \text{ for all } s \leq t,\; \mathcal{W} \models \mathcal{C},\; \mathcal{W} \text{ is consistent with } H_t \right\}.
\]
\end{definition}

The constraint system $\mathcal{C}$ plays an essential role. Without it, $\mathcal{R}$ is excessively large: any world that happens to render to the correct pixel values at all observed times is technically admissible. Physical laws, object permanence, and causal structure eliminate vast subsets of technically pixel-consistent hypotheses.

\begin{remark}
In practice, the constraint system $\mathcal{C}$ will be heterogeneous and only partially formal. Some constraints are physical (rigid bodies do not interpenetrate), some are perceptual (human faces have specific proportions), some are causal (a door that was closed at time $s$ cannot be open at time $s+\epsilon$ without an intervening opening event), and some are user-specified (a specific character has been declared to be wearing a red coat). The framework treats these uniformly as constraints on the admissible set, regardless of their origin.
\end{remark}

\subsection{Admissible Continuations}

The forward problem, equally important for generation, is: given the current state $\mathcal{W}_t$ and its history $H_t$, what future states are admissible?

\begin{definition}[Admissible Continuation Set]
The \textbf{admissible continuation set} at time $t$ is
\[
\mathcal{A}(H_t) = \left\{ \mathcal{W}_{t+1} \;:\; \mathcal{W}_{t+1} \text{ is reachable from } \mathcal{W}_t \text{ by a physically and causally consistent transition, compatible with } \mathcal{C} \text{ and } H_t \right\}.
\]
\end{definition}

The continuation set $\mathcal{A}(H_t)$ is the generative system's operating space. Rather than sampling unconditionally from a learned distribution, a well-formed generative step selects an element of $\mathcal{A}(H_t)$. The admissibility condition enforces coherence.

\begin{remark}
$\mathcal{A}(H_t)$ will typically be very large. A movie frame at time $t$ is consistent with countless physically possible next frames. The narrowing of $\mathcal{A}(H_t)$ by genre, character psychology, dramatic pacing, user instructions, and narrative convention is the subject of Section~\ref{sec:narrative}. The point here is only that $\mathcal{A}(H_t)$ is always a proper subset of all possible world states, bounded by the physical and causal structure committed thus far.
\end{remark}

\subsection{The Probate Mechanism}

A central architectural commitment is that a generated or inferred candidate continuation must not become authoritative merely because a generative model proposed it. This is captured by a probate step, formalized as follows.

\begin{definition}[Probate]
A candidate continuation $\hat{\mathcal{W}}_{t+1}$ proposed by process $P_i$ is said to \textbf{pass probate} if
\[
\hat{\mathcal{W}}_{t+1} \in \mathcal{A}(H_t) \quad \text{and} \quad \Adm(\hat{\mathcal{W}}_{t+1},\; H_t,\; \mathcal{C}) = \textsc{true}.
\]
A candidate that fails probate is not committed to the world state. It may trigger additional witness queries, may be modified by a repair process, or may be rejected in favor of an alternative proposal.
\end{definition}

The sequence proposal $\to$ diagnosis $\to$ admissibility check $\to$ commitment is the governing rhythm of the reconstruction loop. Algorithm~\ref{alg:probate} gives the pseudocode.

\begin{algorithm}[h]
\caption{Probate Loop for World State Continuation}
\label{alg:probate}
\begin{algorithmic}[1]
\Require Current world state $\mathcal{W}_t$, history $H_t$, constraint system $\mathcal{C}$, algorithm registry $\Pi$
\Ensure Updated world state $\mathcal{W}_{t+1}$ and history $H_{t+1}$

\State $\hat{\mathcal{W}} \gets \emptyset$; $\;k \gets 0$
\While{$\hat{\mathcal{W}} \notin \mathcal{A}(H_t)$ \textbf{and} $k < k_{\max}$}
    \State Select proposal process $P_i \in \Pi$ according to scheduler
    \State $\hat{\mathcal{W}} \gets P_i(\mathcal{W}_t,\; H_t)$ \Comment{Proposal}
    \State $D \gets \textsc{Diagnose}(\hat{\mathcal{W}},\; \mathcal{W}_t,\; H_t)$ \Comment{Structural diagnosis}
    \If{$D$ contains critical violations}
        \State Invoke additional witnesses for flagged regions \Comment{See Algorithm~\ref{alg:witness}}
        \State $\hat{\mathcal{W}} \gets \textsc{Repair}(\hat{\mathcal{W}},\; D,\; H_t)$ \Comment{Attempt repair}
    \EndIf
    \State $a \gets \Adm(\hat{\mathcal{W}},\; H_t,\; \mathcal{C})$ \Comment{Admissibility check}
    \State $k \gets k + 1$
\EndWhile
\If{$\hat{\mathcal{W}} \in \mathcal{A}(H_t)$}
    \State $\mathcal{W}_{t+1} \gets \hat{\mathcal{W}}$
    \State $H_{t+1} \gets H_t \cup \{(\hat{\mathcal{W}},\; P_i,\; D,\; a,\; \text{committed})\}$ \Comment{Append provenance}
\Else
    \State $\mathcal{W}_{t+1} \gets \mathcal{W}_t$ \Comment{Hold state; flag for review}
    \State $H_{t+1} \gets H_t \cup \{(\hat{\mathcal{W}},\; P_i,\; D,\; a,\; \text{rejected})\}$
\EndIf
\State \Return $\mathcal{W}_{t+1},\; H_{t+1}$
\end{algorithmic}
\end{algorithm}

\begin{remark}
The repair step in Algorithm~\ref{alg:probate} is intentionally left abstract. A repair may involve invoking a specialized correction model for the failing subregion, reverting to a prior committed state for that region, interpolating from adjacent frames, or substituting an alternative proposal. The framework specifies that repair is attempted before rejection, and that all attempted repairs are recorded in $H_{t+1}$. Specific repair procedures are architectural choices that fall outside the formal framework proper.
\end{remark}

%% ============================================================
\section{Heterogeneous Algorithm Repertoire}
\label{sec:repertoire}
%% ============================================================

\subsection{No Privileged Algorithm}

A significant feature of this framework is the explicit denial of algorithmic privilege. No single algorithm, and in particular no neural generative model, is treated as the authoritative source of world state updates. Instead, the system maintains a \textbf{registry} of procedures, each capable of proposing updates to some component of $\mathcal{W}_t$, and a \textbf{scheduler} that selects and combines those proposals.

This is not merely a design preference but follows from the reconstruction argument. If the goal is admissible reconstruction of a world state consistent with accumulated evidence, then the appropriate procedure is the one best suited to the evidence at hand. Classical multi-view geometry algorithms \cite{hartley2003multiple} may be more reliable than a neural model for estimating depth from stereo pairs with well-characterized cameras. Rigid-body physics simulation is more reliable than a temporal diffusion model for predicting the trajectory of a thrown object under known gravity. A face-tracking system with a calibrated skeleton may be more reliable than a generative model for maintaining identity through rapid head motion.

\begin{proposal}[Algorithm Registry]
\label{prop:registry}
The system maintains a registry $\Pi = \{P_1, P_2, \ldots, P_K\}$ of procedures. Each procedure $P_i$ is characterized by:
\begin{enumerate}[leftmargin=2em,itemsep=0pt]
\item its \textbf{input signature}: which components of $\mathcal{W}_t$ and $H_t$ it reads;
\item its \textbf{output signature}: which components of $\mathcal{W}_t$ it proposes to update;
\item a \textbf{confidence estimator} $\kappa_i : \mathcal{W}_t \times H_t \to [0,1]$ that estimates its expected reliability given the current state;
\item a \textbf{computational cost estimate} $c_i$;
\item an \textbf{applicability predicate} $\alpha_i : \mathcal{W}_t \to \{0,1\}$ indicating whether the procedure's preconditions are met.
\end{enumerate}
\end{proposal}

The registry may include, without any commitment to exhaustiveness: rasterization pipelines, ray and path tracers, voxel-based renderers, Gaussian splatting renderers \cite{kerbl2023gaussians}, neural radiance representations \cite{mildenhall2020nerf}, procedural generation systems, diffusion models, flow-matching models, optical flow estimators \cite{horn1981optical,lucas1981optical}, object trackers, depth estimators, SLAM-like reconstruction pipelines \cite{davison2007monoslam}, inverse kinematics solvers, rigid- and soft-body physics simulators, particle systems, spline and trajectory interpolators, classical signal processing filters, conventional video codecs (used as temporal predictors), deterministic rules and hard constraints, and hand-authored procedural specifications.

\subsection{Scheduling and Selective Invocation}

The governing principle of scheduling is not exhaustive invocation but selective invocation: invoke additional witnesses where uncertainty, disagreement, structural importance, or failure warrants the additional computation. Running every procedure on every frame region at every timestep is not the design goal and is computationally infeasible. The design goal is to use a cheap primary process (potentially a neural temporal model) for the bulk of routine continuations, and to escalate to additional witnesses precisely when and where the primary process shows signs of structural failure.

\begin{proposal}[Witness Scheduler]
The witness scheduler operates in two passes per timestep. In the \textbf{first pass}, a cheap confidence estimator $\kappa_{\text{fast}}$ scans the candidate continuation $\hat{\mathcal{W}}$ and produces a suspicion field $\Sigma : \mathcal{W} \to [0,1]$, where $\Sigma(q)$ is large wherever the fast estimator is uncertain, wherever structural invariants appear to have been violated, or wherever the candidate continuation is far from the prior. In the \textbf{second pass}, for each region $q$ with $\Sigma(q) > \tau$ (where $\tau$ is a configurable threshold), the scheduler selects witnesses from $\Pi$ based on their applicability predicates, confidence estimators, and cost estimates, and invokes them to produce an extended witness set $W(q)$.
\end{proposal}

This hierarchical approach ensures that computational resources scale with difficulty rather than with frame size.

%% ============================================================
\section{Witnesses and Disagreement}
\label{sec:witnesses}
%% ============================================================

\subsection{The Witness Set}

\begin{definition}[Witness Set]
For a query $q$ (which may be a region, an object, a transition, a geometric relationship, or any testable aspect of the world state candidate $\hat{\mathcal{W}}$), let $\mathcal{P}_q \subseteq \Pi$ be the set of applicable procedures at $q$. The \textbf{witness set} is
\[
W(q) = \left\{ P_i(q) \;:\; i \in \mathcal{P}_q \right\}.
\]
Elements of $W(q)$ are candidate values or assessments of the query $q$ produced by independent procedures.
\end{definition}

\begin{remark}
The independence in ``independent procedures'' is meaningful. When optical flow, object tracking, and 3D projection independently predict where a surface point should appear in the next frame, their agreement provides evidence for the prediction while their disagreement provides evidence that something is structurally wrong. This is structurally analogous to ensemble methods in machine learning \cite{lakshminarayanan2017deep}, but applied specifically to world-state properties rather than classification or regression targets.
\end{remark}

\subsection{Dispersion and Disagreement}

\begin{definition}[Dispersion Functional]
The \textbf{dispersion} of the witness set $W(q)$ is a functional
\[
\Gamma(q) = \Disp(W(q)) \in [0, \infty)
\]
quantifying the spread or inconsistency of witness predictions. High $\Gamma(q)$ indicates that independent procedures disagree about the state or behavior of $q$.
\end{definition}

The appropriate choice of dispersion functional depends on the type of $q$:

\begin{itemize}[leftmargin=2em,itemsep=2pt]
\item \textbf{Continuous position or depth predictions.} Let $W(q) = \{v_1, \ldots, v_m\} \subset \mathbb{R}^d$. Then $\Gamma(q) = \frac{1}{m}\sum_{i=1}^m \|v_i - \bar{v}\|^2$ (mean squared deviation from the ensemble mean) or a more robust alternative such as the mean absolute deviation from the median.

\item \textbf{Object identity (categorical).} Let $W(q) = \{c_1, \ldots, c_m\}$ where each $c_i$ is a categorical identity assignment. Then $\Gamma(q) = 1 - \max_c \frac{|\{i : c_i = c\}|}{m}$ (one minus the plurality fraction).

\item \textbf{Optical flow vectors.} Let $W(q) = \{f_1, \ldots, f_m\} \subset \mathbb{R}^2$. Then $\Gamma(q) = \max_{i,j} \|f_i - f_j\|$ (maximum pairwise deviation), which is sensitive to single catastrophically wrong estimates.

\item \textbf{Depth maps.} Disagreement between stereo estimation, monocular estimation, and SLAM-derived depth at pixel $(u,v)$ can be measured as $\Gamma(u,v) = \mathrm{std}(d_{\text{stereo}}(u,v),\; d_{\text{mono}}(u,v),\; d_{\text{SLAM}}(u,v))$.

\item \textbf{Occlusion boundaries.} Binary edge maps from different processes can be compared via a boundary F-measure or a Hausdorff distance between estimated boundary locations.
\end{itemize}

\begin{hypothesis}
\label{hyp:gamma}
The dispersion field $\Gamma$ localizes deliberately introduced structural errors (object identity swaps, impossible geometry, occlusion boundary violations, unphysical motion fields) more reliably than frame-local pixel-domain error metrics computed between consecutive frames.
\end{hypothesis}

\subsection{Disagreement as Diagnostic Information}

The key reframing is: disagreement is information, not merely variance to be minimized or averaged away.

When $\Gamma(q)$ is high, the appropriate response is not to average the witnesses (which can produce a physically impossible composite) but to ask \textit{which witness is more likely to be correct given the available evidence}, and further, \textit{what assumption, representation, or model is responsible for the disagreement}. If optical flow predicts a surface point at position $p_1$, rigid-body projection predicts $p_2$, and object tracking predicts $p_3$, the disagreement is evidence that at least one of the underlying models is failing at this point. Identifying which one fails is the purpose of the diagnostic engine.

\begin{proposal}[Disagreement Response Protocol]
When $\Gamma(q) > \tau_{\text{high}}$, the system should:
\begin{enumerate}[leftmargin=2em,itemsep=0pt]
\item Log $q$ and $W(q)$ to the provenance record $H_t$ with full witness attribution.
\item Attempt to identify the likely outlier witness by checking each $P_i(q)$ against priors, physical constraints, and consistency with neighboring regions.
\item If an outlier is identified, exclude it and recompute the consensus.
\item If no consensus is reachable, flag $q$ as a region of unresolved structural uncertainty, hold the prior state for $q$, and request additional computation or user input.
\item Never silently commit a high-dispersion region as if it were a well-supported update.
\end{enumerate}
\end{proposal}

Algorithm~\ref{alg:witness} gives the pseudocode for witness-based structural diagnosis.

\begin{algorithm}[h]
\caption{Witness-Based Structural Diagnosis}
\label{alg:witness}
\begin{algorithmic}[1]
\Require Candidate continuation $\hat{\mathcal{W}}$, prior state $\mathcal{W}_t$, history $H_t$, registry $\Pi$, thresholds $\tau_{\text{fast}}, \tau_{\text{high}}$
\Ensure Diagnosis record $D$, updated candidate $\hat{\mathcal{W}}'$

\State $\Sigma \gets \kappa_{\text{fast}}(\hat{\mathcal{W}},\; \mathcal{W}_t)$ \Comment{Fast suspicion field over all queries}
\State $D \gets \emptyset$
\For{each query $q$ with $\Sigma(q) > \tau_{\text{fast}}$}
    \State $\mathcal{P}_q \gets \{ P_i \in \Pi : \alpha_i(\mathcal{W}_t) = 1,\; P_i \text{ is applicable to query type of } q\}$
    \State $W(q) \gets \{ P_i(q) : i \in \mathcal{P}_q \}$
    \State $\Gamma(q) \gets \Disp(W(q))$
    \If{$\Gamma(q) > \tau_{\text{high}}$}
        \State Identify likely outlier witnesses in $W(q)$
        \State $D \gets D \cup \{(q,\; W(q),\; \Gamma(q),\; \text{outlier estimate})\}$
        \State Mark $q$ as structurally suspect in $\hat{\mathcal{W}}$
    \ElsIf{$\Gamma(q) > \tau_{\text{fast}}$}
        \State $D \gets D \cup \{(q,\; W(q),\; \Gamma(q),\; \text{minor uncertainty})\}$
    \EndIf
\EndFor
\State $\hat{\mathcal{W}}' \gets \hat{\mathcal{W}}$ with suspect regions flagged or held from prior $\mathcal{W}_t$
\State \Return $D,\; \hat{\mathcal{W}}'$
\end{algorithmic}
\end{algorithm}

%% ============================================================
\section{Structural Glitches as Diagnostic Evidence}
\label{sec:glitches}
%% ============================================================

\subsection{Redefining the Glitch}

Conventional usage treats a visual artifact or glitch as primarily an aesthetic failure: something the viewer can see that should not be visible. This definition conflates the detectability of the artifact with its cause and with the kind of information whose loss produced it. We propose a more informative definition.

\begin{definition}[Structural Glitch]
\label{def:glitch}
A \textbf{structural glitch} at query $q$ and time $t$ is a configuration of $\hat{\mathcal{W}}_{t+1}$ that violates a structural relation that was both true in $\mathcal{W}_t$ (or recoverable from $H_t$) and was, in principle, preservable by the transition process given the information available to it. More precisely, let $\phi$ be a structural predicate (such as ``object $o$ and region $r$ are on the same surface'', ``depth increases monotonically across boundary $b$'', ``face identity is consistent with identity $\text{id}$''). A structural glitch occurs when $\phi(\hat{\mathcal{W}}_{t+1}) = \text{false}$ but $\phi(\mathcal{W}_t) = \text{true}$ and the transition process had sufficient information to preserve $\phi$.
\end{definition}

The phrase ``had sufficient information to preserve $\phi$'' is essential and distinguishes this definition from a simple continuity requirement. A compression artifact that violates a boundary relation is not a structural glitch if the bitrate was genuinely too low to represent that boundary. It is a structural glitch if the codec had access to object segmentation and chose to disregard it.

\subsection{A Taxonomy of Structural Glitches}

The following categories organize structurally meaningful failure modes. Each category corresponds to a class of predicates $\phi$ in Definition~\ref{def:glitch}.

\textit{Object identity violations.} The system assigns a different identity to the same physical object between frames, or merges two distinct objects into one, or splits one object into two, without a causally supported reason. This includes the ``extra finger'' problem in generative models and identity drift in tracking.

\textit{Impossible topology.} The geometry of a proposed update is not homeomorphic to any physically realizable surface: a surface self-intersects, a closed surface suddenly opens, limbs detach and reattach.

\textit{Temporal discontinuity.} A state variable changes discontinuously between adjacent frames by more than can be explained by frame rate, motion blur, or committed causal events.

\textit{Motion field inconsistency.} The optical flow predicted from the 3D geometry and camera motion is inconsistent with the actual pixel displacement field in the rendered output.

\textit{Depth inconsistency.} A foreground object appears at a depth that places it behind a background object at the same image location.

\textit{Occlusion errors.} A surface that should be occluded is rendered as visible, or vice versa, in a way inconsistent with the geometry.

\textit{Foreground/background texture leakage.} Texture from one depth layer bleeds into another, the canonical codec glitch at low bitrates and the canonical generative failure at inpainting boundaries.

\textit{Boundary-crossing artifacts.} A compression or generative block boundary crosses a scene object boundary, producing artifacts that do not respect the separation between surfaces.

\textit{Material inconsistency.} The material of an object changes in a way not supported by lighting changes, causing a surface to appear to change its reflectance or color for no physical reason.

\textit{Lighting inconsistency.} A shadow, specular highlight, or ambient contribution does not match the committed illumination model $L_t$.

\textit{Causal inconsistency.} An event recorded in $S_t$ or $H_t$ has consequences that are violated in the proposed continuation.

\textit{Narrative-state inconsistency.} A semantic fact committed to $S_t$ (a character's physical injury, the location of an object, a relationship between characters) is violated without an intervening causal event.

\begin{proposition}
Each category of structural glitch in the taxonomy corresponds to a structural predicate $\phi$ that is checkable, at least in principle, from the world state $\mathcal{W}_t$ and history $H_t$. Therefore, each category defines a class of diagnostic tests that can be applied to candidate continuations without reference to pixel-level appearance alone.
\end{proposition}

\subsection{Glitches as Evidence}

A structural glitch, once detected, is not merely something to suppress. It is evidence about the failure of a particular representation or process. If a foreground/background texture leak occurs at the boundary between a foreground character and the background wall, it is evidence that whichever process produced that region did not have access to, or did not correctly use, the depth and segmentation information that would have distinguished the two surfaces. The appropriate response is to route subsequent updates for that region to a process with better access to that information, or to invoke the depth estimator and object segmenter as additional witnesses.

\begin{remark}
This is a strictly stronger claim than the claim that glitches indicate ``a bad model was used.'' The diagnostic value of a glitch lies in identifying \textit{which structural relation was violated} and thereby identifying \textit{which representation or module failed to preserve it}. Different glitch categories implicate different modules: texture leakage implicates the depth/segmentation pipeline; identity drift implicates the tracking module; shadow inconsistency implicates the lighting model. This specificity is what makes glitches diagnostically useful rather than merely undesirable.
\end{remark}

%% ============================================================
\section{Structure-Aware Rate-Distortion}
\label{sec:ratedistortion}
%% ============================================================

\subsection{Conventional Rate-Distortion}

Classical rate-distortion theory \cite{berger1971,coverandthomas2006} frames compression as a tradeoff between the number of bits used to represent a signal and a distortion measure between the original and the reconstruction. In video compression, the standard objective is approximately

\[
J_{\mathrm{RD}} = R + \lambda D_{\mathrm{pixel}},
\]

where $R$ is the bit rate and $D_{\mathrm{pixel}}$ is a frame-level distortion measure such as mean squared error or a variant.

Modern video codecs (H.264/AVC \cite{wiegand2003}, H.265/HEVC \cite{sullivan2012}) incorporate significant structure-sensitivity beyond pixel-domain distortion: motion estimation and compensation exploit inter-frame coherence, intra-frame prediction exploits spatial correlation, perceptual quantization adapts quantization to the local masking properties of the visual system, and loop filters smooth block-boundary ringing. These are meaningful structural features, and characterizing existing codecs as purely pixel-domain systems would be inaccurate.

\subsection{The Proposed Extension}

The claim of the present framework is more specific than ``codecs should be more perceptual.'' It is that a persistent world representation exposes semantic and geometric invariants that are not available to a codec operating on compressed frame-level features. A codec processing $H \times W$ blocks of pixels has no stable representation of object identity, scene depth, occlusion relationships, or causal state. When it fails under bitrate pressure, its failures have no reason to respect boundaries between foreground and background objects, because those boundaries are not represented in the form the codec operates on.

We propose a structure-aware rate-distortion objective:

\[
J_{\mathrm{struct}} = R + \lambda_1 D_{\mathrm{appearance}} + \lambda_2 D_{\mathrm{boundary}} + \lambda_3 D_{\mathrm{motion}} + \lambda_4 D_{\mathrm{identity}} + \lambda_5 D_{\mathrm{depth}} + \lambda_6 D_{\mathrm{temporal}},
\]

where each $D_{\mathrm{(\cdot)}}$ is a distortion measure specific to a structural category:

\begin{itemize}[leftmargin=2em,itemsep=2pt]
\item $D_{\mathrm{appearance}}$: a perceptually weighted pixel-level distortion (SSIM \cite{wang2004ssim} or LPIPS), treating appearance as only one component rather than the sole component;
\item $D_{\mathrm{boundary}}$: a boundary-fidelity term measuring how well object and depth boundaries are preserved, computed from the object segmentation and depth map available in $\mathcal{W}_t$;
\item $D_{\mathrm{motion}}$: distortion in the motion field, measured between the 3D-geometry-derived flow and the decoded-frame flow;
\item $D_{\mathrm{identity}}$: a term penalizing object-identity violations, computed from object-level feature embeddings rather than raw pixels;
\item $D_{\mathrm{depth}}$: distortion in the decoded depth map relative to the world-model depth;
\item $D_{\mathrm{temporal}}$: a temporal consistency term penalizing frame-to-frame oscillations in low-frequency regions that should be stable.
\end{itemize}

\begin{remark}
The $\lambda_i$ weights are hyperparameters that depend on application requirements. A system compressing action cinema at very low bitrates may prioritize $\lambda_3$ (motion) and $\lambda_2$ (boundary); a system compressing talking-head video may prioritize $\lambda_4$ (identity). The framework does not specify universal weights but makes explicit the structure of the objective.
\end{remark}

\subsection{The Informational Claim}

The deeper claim motivating the structural objective is informational. A representation that discards distinctions about object identity, depth, and boundaries at an early stage cannot reliably preserve those distinctions at later stages, regardless of the distortion measure used in optimization. This is not a criticism of any specific codec design but a statement about information flow: structure that is not represented cannot be preserved.

\begin{proposition}
\label{prop:infoclaim}
Let $\phi$ be a structural predicate (e.g., ``pixels at positions $p_1$ and $p_2$ belong to different objects''). If the compressed representation $\hat{I}$ does not contain sufficient information to evaluate $\phi$, then $\phi$ cannot be guaranteed to hold in $\hat{I}$ regardless of the decoding procedure used.
\end{proposition}

\begin{proof}[Justification]
This is a direct consequence of the data processing inequality: post-decoding processing cannot recover information that was not retained in the compressed stream. If the encoder discarded the segmentation boundary between two objects, the decoder cannot reconstruct it from pixel data alone in any representation-independent sense.
\end{proof}

\begin{hypothesis}
\label{hyp:structure-ratedistortion}
Under aggressive bitrate constraints, a system encoding video with access to world-state structure (depth, segmentation, object identity) will produce artifacts that are qualitatively different from those produced by a structure-unaware codec: the structured system's artifacts will tend to remain within individual objects or regions rather than crossing object boundaries, even when comparable total distortion is measured by $D_{\mathrm{appearance}}$ alone.
\end{hypothesis}


%% ============================================================
\section{Style as a Realization Operator}
\label{sec:style}
%% ============================================================

\subsection{Separating World Identity from Visual Realization}

Style, in the broad sense employed here, denotes the family of choices that determine the visual character of a rendered output without determining the identity of what is depicted. Two animations of the same scene, one rendered as photorealistic cinema and one rendered in a charcoal-drawing style, depict the same world. The same dog is behind the same table, the same moon is at the same position in the sky, and the same character has the same number of arms.

\begin{definition}[Style Realization Operator]
A \textbf{style realization operator} is a function
\[
S: \mathcal{W} \times \Theta \to V,
\]
where $\mathcal{W}$ is the world state, $\Theta$ is a style parameter space, and $V$ is the space of rendered outputs (images, videos, or other visual formats). Two operators $S_1 = S(\cdot, \theta_1)$ and $S_2 = S(\cdot, \theta_2)$ define the same \textbf{world realization} if $S_1(\mathcal{W}) \neq S_2(\mathcal{W})$ for some $\mathcal{W}$, but they share the same underlying world.
\end{definition}

The same world state thus admits multiple realizations:
\[
\mathcal{W} \;\xrightarrow{S_1}\; V_1, \qquad \mathcal{W} \;\xrightarrow{S_2}\; V_2.
\]

Photorealistic rendering, VHS-era video, charcoal animation, cel animation, pixel art, stop-motion clay aesthetic, abstract motion graphics, blueprint-style line drawing, and typewriter ASCII video are different style realization operators applied to the same world. Their outputs may look radically different while preserving selected geometric, causal, and narrative properties.

\subsection{Distinguishing Invariants from Style Dimensions}

Not every visual property is a style dimension. Some properties are invariants that a style realization must preserve; others are dimensions that style is permitted to transform.

\begin{definition}[Style Invariant]
A property $\phi$ of $\mathcal{W}$ is a \textbf{style invariant} for a family of style operators $\{S_\theta : \theta \in \Theta\}$ if, for all $\theta \in \Theta$ and all $\mathcal{W}$, the property $\phi$ can be recovered from $S_\theta(\mathcal{W})$.
\end{definition}

Style invariants include, at minimum: object count and approximate spatial layout, depth ordering between foreground and background objects, motion direction (an object moving left continues to move left across style variants), occlusion relationships, and narrative-level facts recorded in $S_t$ (a character who has been wounded is wounded in every style variant of the same world state).

Style dimensions, by contrast, include: color palette, spatial resolution and pixelation, edge rendering (sharp, soft, sketched, nonexistent), texture fidelity and material appearance, lighting softness and color temperature, temporal frame rate and motion blur character, and grain or noise character.

\begin{proposition}
If a style realization operator $S_\theta$ is applied to a world state $\mathcal{W}$ and a subsequent style realization operator $S_{\theta'}$ is applied to the same $\mathcal{W}$, the difference $S_\theta(\mathcal{W}) - S_{\theta'}(\mathcal{W})$ (in any appropriate metric) should contain no information about world-level properties that are style invariants.
\end{proposition}

This proposition defines a consistency requirement on the style system: changing style should not change the world. It provides a testable condition for style-transfer systems that operate on the world state rather than on frames.

\begin{remark}
Existing style-transfer systems that operate on frames or on latent frame representations do not generally satisfy this condition, because frame-level stylization can alter structural properties as a side effect. The world-state approach provides an architectural guarantee: because style is applied after world-state update, style choices cannot retroactively alter the world's geometric or causal structure.
\end{remark}

%% ============================================================
\section{Narrative Continuation and Story State}
\label{sec:narrative}
%% ============================================================

\subsection{Extending Persistence to Semantic Facts}

The world state $\mathcal{W}_t$ includes a semantic and narrative component $S_t$. This component is responsible for the persistence of facts that are not directly recoverable from geometry alone: a character's knowledge of another character's secret, the unfired gun introduced in the first act, the promise made between two characters, the physical injury a character sustained three scenes earlier.

\begin{definition}[Story State]
The \textbf{story state} $S_t \subseteq \mathcal{F}$ is a set of semantic facts indexed by time, where $\mathcal{F}$ is a universe of representable story facts. A fact $f \in S_t$ may be geometrically grounded (``character $A$ is in room $R$''), relational (``character $A$ knows that character $B$ is the murderer''), causal (``the bridge was destroyed at time $s < t$''), possessive (``character $A$ holds object $O$''), or evaluative (``the dramatic tension of the scene is high'').
\end{definition}

The story state is cumulative: unless a fact is explicitly negated by a committed causal event, it persists. A character who is injured at time $s$ is injured at all subsequent times $t > s$ unless a healing event has been committed to $H_t$. This persistence is what distinguishes the narrative engine from a frame-local language model prompted with a summary of previous events.

\subsection{Admissible Narrative Continuation}

\begin{definition}[Narrative Admissibility]
A proposed story-state transition $S_{t+1}$ is \textbf{narratively admissible} from $S_t$ if it is obtainable from $S_t$ by: (a) committing facts that follow from physical causality given committed geometric events, (b) committing facts that follow from causal events explicitly authored or proposed, or (c) leaving facts unchanged. It is \textbf{narratively inadmissible} if it negates a persistent fact without a committed causal explanation, or if it introduces a fact that contradicts a committed fact without explicit override.
\end{definition}

\begin{proposal}
Plot generation operates as constraint-narrowing over the admissible narrative continuation set $\mathcal{A}(H_t) \cap \mathcal{A}_{\mathrm{narrative}}(S_t)$. Genre, character psychology, physical law, pacing, dramatic structure, and user instructions are each formalized as constraints that progressively reduce this set. A detective story and an absurdist comedy can share the same world engine and the same geometry while possessing radically different continuation operators, because their genre constraints narrow $\mathcal{A}$ in different directions.
\end{proposal}

\begin{remark}
Intentional violations of narrative admissibility (surrealist discontinuity, unreliable narration, dream sequences with impossible geometry, fourth-wall breaks) should be represented as changes to the governing constraint system $\mathcal{C}$ rather than as accidental failures of the narrative engine. A dream sequence is not a narrative glitch; it is the activation of a different constraint regime. The system should mark such transitions explicitly in $H_t$.
\end{remark}

%% ============================================================
\section{Camera as Epistemic Operator}
\label{sec:camera}
%% ============================================================

\subsection{The Camera Does Not Create the World}

The observation map $\Obs_C : \mathcal{W} \to I$ reveals a partial projection of the world rather than creating it. Moving the camera changes $\Obs_C$ but does not, in general, change $\mathcal{W}$. This is the formal counterpart of the cinematographer's intuition: framing is selection, not invention.

\begin{definition}[Cinematographic Operations as Observation Changes]
The following operations modify the observation operator without modifying the underlying world state:
\begin{enumerate}[leftmargin=2em,itemsep=0pt]
\item \textbf{Pan}: rotation of the camera about a vertical or horizontal axis, changing the visible frustum;
\item \textbf{Tilt}: rotation about the camera's lateral axis;
\item \textbf{Zoom}: change in field of view (focal length), equivalent to a change in the projection parameters of $C_t$;
\item \textbf{Dolly/track}: translation of the camera in the world, changing the center of projection;
\item \textbf{Rack focus}: change in focus depth, altering the depth-of-field parameters of the rendering;
\item \textbf{Cut}: discontinuous jump to a new camera configuration $C'$, equivalent to switching observation operators between frames;
\item \textbf{POV shot}: assigning $C_t$ to a camera position and orientation coinciding with a character's viewpoint, as specified in $O_t$.
\end{enumerate}
\end{definition}

\begin{remark}
This formalization connects cinematography to the theory of partially observable systems. A cut is a discontinuous observation operator change. A POV shot defines an observation operator that is a function of a character's state. A surveillance-camera shot is an observation operator fixed in world coordinates regardless of character positions. Each of these is formally distinct, and their distinctness can be exploited for editing: the same world state $\mathcal{W}_t$ yields different frames under each, without requiring re-generation of the world itself.
\end{remark}

\subsection{Hidden State and Off-Screen Persistence}

Because the world state $\mathcal{W}_t$ is maintained independently of what is currently observed, the framework naturally represents hidden state: objects, characters, and events that are present in the world but not currently visible through any active camera.

This has several practical consequences. First, objects off-screen are not removed from $O_t$ but simply not projected by the current $\Obs_{C_t}$. Their dynamics continue: a character off-screen continues to age, move, act, and accrue causal consequences. Second, when the camera returns to observe a previously off-screen region, the world state at that region is determined by its committed continuation, not by a fresh generative sample. Third, the distinction between world-level facts and camera-established facts (facts that exist only because a particular camera observed them at a particular time) becomes formally representable in $H_t$.

\begin{hypothesis}
\label{hyp:offscreen}
A world-state-persistent system maintains more consistent spatial and causal relationships in regions that leave and re-enter frame during a scene than a frame-local temporal model conditioned on a fixed context window, when measured on scripted scenes with known ground truth.
\end{hypothesis}

%% ============================================================
\section{Scene Architecture}
\label{sec:scene}
%% ============================================================

\subsection{Heterogeneous Scene Representation}

The persistent world state $\mathcal{W}_t$ must be representable in a form that supports heterogeneous geometric representations simultaneously, rather than committing to a single universal format such as polygon meshes.

\begin{proposal}[Heterogeneous Scene Graph]
The scene is organized as a directed acyclic graph (the scene graph) where each node $n$ carries:
\begin{enumerate}[leftmargin=2em,itemsep=0pt]
\item An \textbf{object identity} $\mathrm{id}(n)$: a stable unique identifier that persists regardless of changes to the node's geometric representation;
\item A \textbf{geometric representation} $\mathrm{geo}(n)$, which may be a polygon mesh, a voxel grid, a signed-distance field, a Gaussian splat cluster, a neural radiance cache, a procedural generator, a skeletal rig with attached mesh, a particle system, a volumetric medium, a curve or spline, or a purely semantic placeholder with no geometric realization;
\item A \textbf{material specification} $\mathrm{mat}(n)$, which may be a BSDF, a procedural material, a neural material representation, or a style annotation;
\item A \textbf{transform} $T(n) \in SE(3)$: the object's pose in world coordinates;
\item A \textbf{provenance record} $\mathrm{prov}(n)$: which process proposed this node, when, and with what supporting evidence;
\item An \textbf{uncertainty distribution} $\mathrm{unc}(n)$: a distribution over poses, shapes, or identities, representing unresolved ambiguity about this node;
\item A set of \textbf{semantic annotations} $\mathrm{sem}(n)$: category labels, character names, story-state references, and other structured metadata.
\end{enumerate}
Edges in the scene graph encode parent-child relationships (for hierarchical objects), occlusion relationships (from committed depth estimates), attachment relationships (one object held by another), and reference relationships (a character's line-of-sight pointing to an object in $O_t$).
\end{proposal}

\subsection{Object Identity Independent of Representation}

Object identity $\mathrm{id}(n)$ is explicitly decoupled from representation $\mathrm{geo}(n)$. The same object may be represented as a Gaussian splat at one timestep, as a mesh at another (if a reconstruction process produces a mesh from new views), and as a neural radiance field at a third. The object identity and its accumulated causal and narrative history persist across these representation changes.

This independence is crucial for the witness mechanism: two different geometric representations of the same object are two witnesses to the same ground truth, not two different objects.

\subsection{Proposed Data Schema}

Table~\ref{tab:schema} provides a simplified data schema for a scene graph node. Serialization is proposed in a format that supports append-only provenance records, heterogeneous geometric payloads, and streaming updates from multiple concurrent processes.

\begin{table}[h]
\centering
\caption{Simplified scene graph node schema.}
\label{tab:schema}
\begin{tabular}{lll}
\toprule
Field & Type & Description \\
\midrule
\texttt{id} & UUID & Stable object identity \\
\texttt{geo\_type} & enum & Representation type (mesh, splat, SDF, neural, ...) \\
\texttt{geo\_payload} & bytes & Serialized geometric data \\
\texttt{mat\_id} & UUID & Reference to material record \\
\texttt{transform} & float[16] & World-to-object transform (row-major $4\times4$) \\
\texttt{uncertainty} & blob & Distribution over pose/shape, format-dependent \\
\texttt{sem\_labels} & string[] & Semantic annotations \\
\texttt{prov\_entries} & record[] & Append-only provenance log (see Section~\ref{sec:provenance}) \\
\texttt{parent\_id} & UUID? & Parent node in scene graph, if any \\
\texttt{created\_at} & int & Logical timestep of creation \\
\texttt{updated\_at} & int & Logical timestep of last update \\
\bottomrule
\end{tabular}
\end{table}

%% ============================================================
\section{Provenance and Evidentiary History}
\label{sec:provenance}
%% ============================================================

\subsection{The Evidentiary Requirement}

Every committed state transition in the world state should retain sufficient provenance to answer, at any later time:

\begin{enumerate}[leftmargin=2em,itemsep=0pt]
\item What process proposed this state or transition?
\item What observations or evidence supported the proposal?
\item Which constraints were checked, and did they pass?
\item Which witnesses were consulted, and did they agree?
\item What was uncertain at the time of commitment?
\item Was any repair applied before commitment, and if so, by what process?
\item Has this committed state been subsequently revised, and if so, why?
\end{enumerate}

This is an evidentiary structure, not merely a version log. The distinction matters: a version log records what changed, while an evidentiary structure records why a particular state was believed to be correct at the time it was committed.

\subsection{Append-Oriented History}

\begin{proposal}[Provenance Record]
Each committed update to a scene graph node $n$ appends a provenance entry of the form:
\[
\mathrm{prov\_entry} = (\mathrm{timestep},\; \mathrm{proposer},\; \mathrm{evidence},\; \mathrm{witnesses},\; \Gamma,\; \mathrm{constraints\_checked},\; \mathrm{repair\_applied},\; \mathrm{outcome}).
\]
Here $\mathrm{proposer}$ is the identifier of the process $P_i$ that proposed the update; $\mathrm{evidence}$ is a reference to the observations and world-state elements used; $\mathrm{witnesses}$ is the set of witness processes consulted and their outputs; $\Gamma$ is the measured dispersion; $\mathrm{constraints\_checked}$ lists the structural predicates evaluated; $\mathrm{repair\_applied}$ records any correction applied; and $\mathrm{outcome}$ is one of \textsc{committed}, \textsc{repaired\_and\_committed}, or \textsc{rejected}.
\end{proposal}

History is append-oriented: entries are added but previous entries are never overwritten. A revision does not delete the prior committed state; it appends a new entry with $\mathrm{outcome} = \textsc{revised}$ and a reference to the revising entry. This structure supports retrospective analysis of failure chains: if a glitch is detected at time $t$, the provenance record can be traversed to identify the earliest entry in the chain that should have been rejected.

\begin{remark}
Three distinct histories are maintained. The \textbf{world history} records committed updates to $\mathcal{W}_t$. The \textbf{observation history} records the sequence of frames $(I_t)$ that have been rendered from $\mathcal{W}_t$, along with the observation operators used. The \textbf{edit history} records user-initiated interventions: explicit changes to world state, style switches, narrative overrides, and corrections to detected glitches. The three histories are linked but distinct.
\end{remark}


%% ============================================================
\section{Modular System Architecture}
\label{sec:architecture}
%% ============================================================

\subsection{Module Overview}

The architecture is organized as a set of modules with defined contracts. The central invariant is that \textbf{only the Commitment Engine may write to the World State Store}; all other modules produce proposals, observations, or assessments that are routed through the Commitment Engine via the Admissibility Engine. This invariant prevents any single module from unilaterally altering the persistent world state.

\begin{description}[leftmargin=2em,itemsep=6pt]

\item[\textbf{World State Store (WSS).}] The authoritative persistent representation of $\mathcal{W}_t$. Provides read access to all modules; accepts writes exclusively from the Commitment Engine. Maintains the scene graph, material library, illumination state, camera state, story state, and history. Supports queries by object identity, spatial region, temporal range, and semantic category.

\item[\textbf{History/Provenance Store (HPS).}] Append-only log of all provenance entries. Never overwritten; new entries may reference prior entries as superseded. Exposes query APIs for retrieving the full provenance chain of any current world state element.

\item[\textbf{Algorithm Registry (AR).}] Catalog of available procedures $\Pi = \{P_i\}$ with their input/output signatures, confidence estimators $\kappa_i$, cost estimates $c_i$, and applicability predicates $\alpha_i$. The registry is read by the Witness Scheduler to select appropriate witnesses.

\item[\textbf{Proposal Engine (PE).}] Queries the Algorithm Registry and current world state to select a primary proposal process for each timestep. Invokes $P_{\mathrm{primary}}(\mathcal{W}_t, H_t)$ and forwards the candidate $\hat{\mathcal{W}}$ to the Diagnostic Engine. Does not write to the World State Store.

\item[\textbf{Witness Scheduler (WS).}] Receives the suspicion field $\Sigma$ from the Diagnostic Engine and selects additional witnesses from the Algorithm Registry for regions with $\Sigma(q) > \tau_{\mathrm{fast}}$. Schedules witness invocations subject to cost constraints. Forwards witness outputs to the Diagnostic Engine.

\item[\textbf{Diagnostic Engine (DE).}] Computes the fast suspicion field $\Sigma$ over the candidate $\hat{\mathcal{W}}$, incorporates witness outputs to compute dispersion $\Gamma$, checks structural predicates $\phi$ for each flagged region, and produces a diagnosis record $D$. Forwards $D$ and the flagged candidate to the Admissibility Engine.

\item[\textbf{Admissibility Engine (AE).}] Evaluates $\Adm(\hat{\mathcal{W}}, H_t, \mathcal{C})$ for the candidate continuation. Identifies which constraints in $\mathcal{C}$ are violated, if any. Forwards the admissibility verdict to the Commitment Engine. Does not attempt repair; does not write to the World State Store.

\item[\textbf{Commitment Engine (CE).}] The sole authority for writing to the World State Store. Accepts candidates that pass admissibility, applies any final adjustments authorized by the repair protocol, writes the new world state, and appends provenance entries to the History Store. Rejects inadmissible candidates and triggers re-proposal or state hold.

\item[\textbf{Physics/Constraint Engine (PCE).}] Maintains and applies physical laws and hard constraints (collision, gravity, fluid dynamics, rigid-body contact, joint limits) as components of $\mathcal{C}$. Provides constraint-checking services to the Admissibility Engine and simulation services to the Algorithm Registry as one of its registered processes.

\item[\textbf{Narrative State Engine (NSE).}] Maintains $S_t$, evaluates narrative admissibility for proposed $S_{t+1}$, and provides the story-level constraints on $\mathcal{A}(H_t)$. Accepts genre specifications, character behavior profiles, dramatic pacing constraints, and user-authored facts.

\item[\textbf{Camera/Observation Engine (COE).}] Manages $C_t$ and computes the observation map $\Obs_{C_t}$ for rendering. Supports scripted camera paths, procedurally generated camera motion, user-controlled camera, and character-viewpoint cameras. Provides frame-to-world unprojection services used by depth and reconstruction processes.

\item[\textbf{Style/Realization Engine (SRE).}] Applies the style operator $S_\theta$ to produce the rendered output $V_t = S_\theta(\mathcal{W}_t)$ from the committed world state. May delegate to the Algorithm Registry for specific rendering algorithms. Does not modify the World State Store; operates on a read-only copy of $\mathcal{W}_t$.

\item[\textbf{Structural Analyzer (SA).}] Post-render analysis module that applies structural predicates to the rendered output $V_t$ and the world state $\mathcal{W}_t$ simultaneously, detecting glitches that survive rendering. Feeds back to the Diagnostic Engine. This module enables the diagnostic feedback loop: the rendered video is treated as another observation and subjected to structural analysis.

\item[\textbf{Encoder/Codec Interface (ECI).}] Wraps conventional video encoders and exposes structural metadata (depth maps, segmentation, motion fields, object identities) as auxiliary streams that can inform structure-aware rate-distortion optimization. Accepts $\lambda_i$ weights for the structural distortion objective.

\item[\textbf{Steganographic Selector Layer (SSL).}] Optional module, active only when hyperpleonastic embedding is enabled. Intercepts rendering decisions with admissible alternatives and selects among them according to payload bits. Operates strictly within the admissibility envelope. Described further in Appendix~\ref{app:steg}.

\item[\textbf{Evaluation Harness (EH).}] Manages ground-truth scene databases, applies frame-local and structural metrics, computes dispersion localizations, and generates experiment reports. Used for development and validation only; not active in production pipelines.

\end{description}

\subsection{Module Contracts}

The contract between modules is expressed in terms of reads and writes to the shared World State Store:

\begin{enumerate}[leftmargin=2em,itemsep=0pt]
\item A module may read from the WSS at any time.
\item A module may write to the WSS only if it is the Commitment Engine, or if it is explicitly delegated by the Commitment Engine for an authorized repair operation.
\item A module that produces a proposal or assessment must not cache it as ground truth; all proposals are provisional until committed.
\item A module that fails (crashes, times out, or returns an inadmissible result) must not leave a partial write in the WSS; all writes are transactional.
\item No module may delete from the HPS; the provenance record is immutable once written.
\end{enumerate}

%% ============================================================
\section{Experimental Program}
\label{sec:experiments}
%% ============================================================

\subsection{Design Principles}

All experiments begin with controlled synthetic scenes for which the complete latent world state $\mathcal{W}_t$ is known. This ground-truth availability is what separates the proposed experiments from evaluations of generative video quality, which typically rely on human preference ratings or reference-based metrics without knowledge of the true scene. Starting from known ground truth allows the measurement of structural errors directly, not merely their downstream visual consequences.

\subsection{Experiment 1: Dispersion Localization of Structural Errors}

\textit{Goal}: Test Hypothesis~\ref{hyp:gamma}, that $\Gamma$ localizes structural errors more reliably than frame-local metrics.

\textit{Setup}: Render a synthetic scene (a room with several objects of known geometry, depth ordering, and identity) from a moving camera. Introduce deliberate structural perturbations: an object identity swap (two objects exchange identity labels at a specified frame), an impossible depth inversion (a foreground object is rendered with the depth of a background object), and a boundary-crossing artifact (a texture block from object $A$ is substituted into a region of object $B$). Each perturbation is introduced at a known time and location.

\textit{Metrics}: For each perturbation, compute (a) the frame-level SSIM and LPIPS between perturbed and unperturbed renderings, (b) the dispersion $\Gamma(q)$ at the perturbation location using at least three witness processes (optical flow, object tracker, depth estimator), and (c) the precision and recall of the dispersion field $\Gamma$ in localizing the perturbation region versus a thresholded SSIM map.

\textit{Success criterion}: $\Gamma$ achieves higher precision and recall at lower thresholds than the SSIM-based detector for each perturbation category.

\subsection{Experiment 2: Object Persistence Through Occlusion}

\textit{Goal}: Test Hypothesis~\ref{hyp:occlusion}.

\textit{Setup}: Render a scene in which a moving object (with known identity, trajectory, and geometry) passes behind a large occluder and reappears on the other side. Compare (a) a world-state-persistent system that maintains the object in $O_t$ through occlusion with (b) a frame-local baseline (a recent-context video diffusion model) conditioned on the last $k$ observed frames. Measure the photometric, geometric, and identity consistency of the reappearing object.

\textit{Metrics}: Identity consistency (whether a face or texture recognition system assigns the same identity before and after occlusion), geometric consistency (whether the 3D bounding box on reappearance is consistent with the pre-occlusion trajectory), and photometric consistency (L2 distance in feature space between the object's appearance before and after occlusion).

\textit{Success criterion}: The world-state-persistent system outperforms the frame-local baseline on all three metrics when the occlusion duration exceeds one second of video at 24fps.

\subsection{Experiment 3: Hierarchical Witness Scheduling Efficiency}

\textit{Goal}: Measure whether selective witness invocation achieves comparable diagnostic coverage to exhaustive invocation at substantially lower computational cost.

\textit{Setup}: Use the synthetic perturbed scenes from Experiment 1. Run (a) exhaustive invocation of all applicable witnesses for all queries at all frames, and (b) hierarchical scheduling using the fast suspicion field $\Sigma$ as a router with threshold $\tau_{\mathrm{fast}}$.

\textit{Metrics}: Total compute time, precision and recall of perturbation localization for both strategies, and the fraction of the total frame area that triggers full witness invocation under the hierarchical strategy.

\textit{Success criterion}: The hierarchical strategy achieves at least 90\% of the precision and recall of exhaustive invocation while invoking full witnesses on at most 20\% of the frame area per timestep.

\subsection{Experiment 4: Structure-Aware Compression}

\textit{Goal}: Test Hypothesis~\ref{hyp:structure-ratedistortion}.

\textit{Setup}: Encode synthetic scenes at varying bitrates using (a) a standard codec (H.265 with default settings) and (b) the proposed Encoder/Codec Interface with structural metadata (depth map, segmentation, object identities) provided as auxiliary streams and structural distortion terms $D_{\mathrm{boundary}}$, $D_{\mathrm{motion}}$, $D_{\mathrm{identity}}$ included in the optimization objective.

\textit{Metrics}: At each bitrate: SSIM and LPIPS (appearance metrics), boundary-crossing artifact rate (fraction of compression block boundaries that cross a known scene object boundary), object identity fidelity (fraction of frames in which the top-1 identity assignment is correct), and the depth consistency error between the decoded video and the ground-truth depth map.

\textit{Success criterion}: At the same SSIM, the structured codec produces fewer boundary-crossing artifacts and higher object identity fidelity than the standard codec.

\subsection{Experiment 5: Style Invariant Preservation}

\textit{Goal}: Test the style realization proposal; specifically, that world-level structural invariants are preserved across style changes.

\textit{Setup}: Render the same world state $\mathcal{W}$ under four style operators: photorealistic, cel animation, charcoal sketch, and pixel art. For each pair of style renderings, evaluate (a) object count and approximate spatial layout consistency, (b) depth ordering consistency, and (c) motion direction consistency for a moving object.

\textit{Metrics}: For each property category, a binary consistency score per frame pair, and an average consistency score across the full scene.

\textit{Success criterion}: Consistency scores exceed 95\% for object count and spatial layout, 90\% for depth ordering, and 95\% for motion direction across all style pairs.

%% ============================================================
\section{Implementation Roadmap}
\label{sec:implementation}
%% ============================================================

\subsection{Staging Principle}

The implementation is staged so that each stage tests a narrow, falsifiable proposition before the next is built. The stages are cumulative: each adds a module or property to the system, using the preceding stages as the tested foundation.

\textit{Stage 0: Baseline Scene and Rendering Infrastructure.} Construct a small controllable 3D scene in Blender \cite{blender2024}. The scene has between five and ten objects of known geometry, depth ordering, identity, and material. Implement a camera trajectory with at least one occlusion event. Render using standard Blender tools and export frames as PNG sequences. Export ground-truth depth maps, segmentation masks, and optical flow fields as reference data. This stage establishes the evaluation harness (corresponding to the Evaluation Harness module) and the Camera/Observation Engine in minimal form.

\textit{Stage 1: World State Representation and Persistent Object Identity.} Implement the World State Store and History/Provenance Store in minimal form, using a flat database of scene graph nodes (Python with SQLite or similar). Manually author the ground-truth world state for the Stage 0 scene. Verify that object identity, depth ordering, and occlusion relationships can be queried correctly at each timestep. This stage tests the data model and establishes the scene architecture independently of any generative or reconstruction process.

\textit{Stage 2: Single-Process Proposal and Admissibility.} Add the Proposal Engine with a single proposal process (a temporal interpolation of object poses between authored keyframes). Add the Admissibility Engine with a single hard constraint (depth ordering: foreground objects may not appear behind background objects). Introduce deliberate pose perturbations that violate depth ordering and verify that the Admissibility Engine rejects them. This stage tests the probate loop (Algorithm~\ref{alg:probate}) in a minimal form.

\textit{Stage 3: Multi-Witness Structural Diagnosis.} Add the Algorithm Registry, Witness Scheduler, and Diagnostic Engine. Register at minimum three witnesses: an optical flow estimator (using OpenCV's Farneback flow \cite{farneback2003}), a depth estimator (MiDaS or similar single-image estimator \cite{ranftl2020}), and the ground-truth depth from the scene. Compute $\Gamma$ over the perturbed scenes from Experiment 1 and verify localization. This is the core of Experiment 1.

\textit{Stage 4: Occlusion Persistence.} Extend the World State Store to maintain object state through occlusion intervals. Implement Experiment 2 and measure reappearance consistency. Compare to the frame-local baseline (any off-the-shelf video temporal model or interpolation).

\textit{Stage 5: Style Realization.} Add the Style/Realization Engine. Implement at least two style operators: the existing Blender rendering (photorealistic) and a post-processing edge-detection and palette-reduction pipeline approximating a sketch style. Implement Experiment 5 and verify invariant preservation.

\textit{Stage 6: Structural Analysis Feedback.} Add the Structural Analyzer as a post-render module. Route its output back to the Diagnostic Engine and trigger re-rendering of flagged regions under alternative style or algorithm choices. This stage completes the diagnostic feedback loop.

\textit{Stage 7: Codec Interface and Compression Experiments.} Wrap FFmpeg \cite{ffmpeg2024} as the Encoder/Codec Interface. Inject depth and segmentation maps as metadata. Implement the structural distortion terms $D_{\mathrm{boundary}}$ and $D_{\mathrm{identity}}$ and run Experiment 4.

The full system is achievable with Python, Blender, FFmpeg, OpenCV, and any off-the-shelf depth estimation and optical flow library. No proprietary or neural-generative components are required through Stage 4. Neural components can be added to the Algorithm Registry in any stage as additional witnesses without disrupting the architecture.

%% ============================================================
\section{Limitations}
\label{sec:limitations}
%% ============================================================

Several limitations of the present framework should be stated clearly, as they determine where further work is required before the framework's empirical hypotheses can be evaluated.

\textit{Computational cost of world state maintenance.} Maintaining a full persistent world state $\mathcal{W}_t$ is substantially more expensive than maintaining a frame buffer. Storing the geometry, materials, object identities, story state, and full provenance history of a complex dynamic scene is a significant memory and compute burden. The hierarchical witness scheduling proposal (Proposal~\ref{prop:registry}) is designed to mitigate this, but the base cost of world-state maintenance has not been empirically characterized for any realistic scene complexity. Whether the diagnostic benefits justify this cost depends on application requirements.

\textit{World state completeness.} The world state $\mathcal{W}_t$ is a model, not reality. For video captured from the real world (rather than synthetically generated), the world state must be estimated from observations, which returns the system to the reconstruction problem. The quality of all downstream processing depends on the quality of this initial reconstruction, and reconstruction from monocular video remains an open research problem.

\textit{Semantic ambiguity in story state.} The formalization of $S_t$ as a set of semantic facts leaves open the question of how to represent graded or uncertain facts (``character $A$ may know character $B$'s secret''), temporally extended facts (ongoing events), and facts that resist propositional encoding (mood, dramatic irony, aesthetic tone). A complete implementation of the Narrative State Engine would require choices about the knowledge representation formalism that the present framework deliberately leaves open.

\textit{Constraint specification burden.} The admissibility framework is only as good as the constraint system $\mathcal{C}$. Specifying complete and consistent physical, causal, and narrative constraints for an arbitrary story in an arbitrary world is a substantial authoring burden. The framework assumes this constraint specification is available but does not address its construction.

\textit{Algorithm repertoire coverage.} The heterogeneous repertoire is only as useful as its coverage of the failure modes that arise in practice. A structural failure that no registered witness is capable of detecting will be committed without diagnosis. Characterizing the coverage of a given registry requires empirical evaluation on a wide range of scene types.

\textit{Distinction between structural and perceptual tolerability.} Not every structural glitch, as defined in Definition~\ref{def:glitch}, is visually perceptible, and not every perceptible artifact corresponds to a structural violation in the sense of the definition. The relationship between the structural and perceptual concepts requires empirical study, likely with human evaluation studies comparing structural error rates to viewer-perceived quality.

%% ============================================================
\section{Conclusion}
\label{sec:conclusion}
%% ============================================================

We have developed a formal framework for video generation, reconstruction, rendering, compression, and diagnosis organized around persistent world state rather than independent frame synthesis. The central commitment is that a rendered frame is an observation of a world, not the world itself. This commitment resolves several structural problems in frame-local approaches: occluded objects persist in the world state, camera motion is an observation operator change rather than a generative challenge, style is a realization operator that leaves world identity invariant, and narrative facts accumulate in a persistent story state that constrains all future continuations.

The framework introduces a heterogeneous algorithm repertoire without algorithmic privilege, a formal witness and disagreement mechanism that converts model disagreement into diagnostic information, and a probate protocol ensuring that candidate continuations must pass structural admissibility checking before commitment to persistent world state. Structural glitches are defined as violations of preserved or recoverable relations, enabling them to serve as diagnostic evidence about which process or representation failed.

A structure-aware rate-distortion objective is proposed that decomposes distortion into semantically meaningful components (appearance, boundary, motion, identity, depth, temporal), extending the informational claim that representations discarding structural distinctions cannot reliably preserve them downstream.

The proposed modular architecture isolates all world-state writes in a single Commitment Engine, routes all proposals through a diagnosis and admissibility pipeline, and maintains an append-only evidentiary history. The architecture is implementable with standard tools from Stage 0 through Stage 4 without any generative neural components, preserving the testability of each stage independently.

The framework's relationship to existing work is one of organization rather than invention of components. Persistent world representations, ensemble disagreement, multi-component distortion objectives, and steganographic embedding are each individually established. The distinctive contribution is their integration under a single admissible reconstruction architecture in which world state is primary, video is observation, and every transformation is assessed against structural invariants that the available information was sufficient to preserve.

%% ============================================================
%% APPENDIX
%% ============================================================

\appendix

\section{Hyperpleonastic Steganography}
\label{app:steg}

\subsection{Overview and Motivation}

The rendering pipeline must make countless decisions that are not fully determined by the semantic and geometric content of $\mathcal{W}_t$. Given a surface that must be rendered as ``rough concrete,'' there may be dozens of procedurally generated microgeometry realizations that are perceptually indistinguishable at the relevant viewing conditions. Given a tree canopy, the precise placement of individual leaves within their admissible spatial envelope may have no effect on the perceived appearance of the canopy. Given a particle system for smoke, the specific initial positions of particles within a volume may produce perceptually equivalent smoke.

Each such decision point is an instance of underdetermination: the rendering constraints do not select a unique realization, and the choice among realizations is in practice arbitrary, random, or driven by aesthetic convenience. The observation is that this freedom is not meaningless; it is a selectable capacity.

\begin{definition}[Admissibility Envelope]
For a scene element $x$ and a set of rendering constraints $\mathcal{C}_x$, the \textbf{admissibility envelope} is the set of realizations consistent with all constraints:
\[
\mathcal{A}(x) = \{ r \;:\; r \text{ satisfies } \mathcal{C}_x \}.
\]
An admissibly equivalent realization is any $r' \in \mathcal{A}(x)$ with $r' \neq r$.
\end{definition}

\subsection{The Steganographic Selector}

\begin{definition}[Hyperpleonastic Steganographic Encoder]
\label{def:steg}
Given a scene element $x$, a payload bit string $b \in \{0,1\}^*$, and an admissibility envelope $\mathcal{A}(x)$ containing $N$ selectable realizations $\{r_0, r_1, \ldots, r_{N-1}\}$, the \textbf{hyperpleonastic steganographic encoder} selects
\[
E(x, b, \mathcal{A}) = r_{f(b)},
\]
where $f : \{0,1\}^{\lfloor \log_2 N \rfloor} \to \{0, 1, \ldots, N-1\}$ is a bijection and $b$ supplies the $\lfloor \log_2 N \rfloor$ bits needed to index the selection. The encoder satisfies the admissibility constraint $E(x, b, \mathcal{A}) \in \mathcal{A}(x)$ for every $b$ by construction, since every $r_i \in \mathcal{A}(x)$.
\end{definition}

The idealized capacity at decision point $x$ is $\lfloor \log_2 N \rfloor$ bits. This is a selection capacity, not a channel capacity in the Shannon sense: it assumes perfect extraction and no transmission noise. Robust recoverable capacity under transcoding, resizing, re-rendering, noise, and color transformation will be substantially lower and requires separate empirical characterization.

\subsection{Binary Selectors and Generalization}

The simplest case is a binary selector: a rendering parameter $p \in [p_{\min}, p_{\max}]$ that is perceptually negligible within the interval is partitioned into a lower basin $B_0 = [p_{\min}, p_{\mathrm{mid}})$ and an upper basin $B_1 = [p_{\mathrm{mid}}, p_{\max}]$. The selector chooses $p \in B_0$ to encode a $0$ bit and $p \in B_1$ to encode a $1$ bit. The rendering decision is indistinguishable at the viewing conditions regardless of which basin is used, but the choice is recoverable by a reader with knowledge of the encoding scheme and access to $\mathcal{A}(x)$.

More generally, for an $n$-ary selector with $n$ admissibly equivalent realizations, the selection capacity is $\lfloor \log_2 n \rfloor$ bits per decision point. Summing over all decision points in a frame, and then over all frames in a video, the aggregate selectable capacity can be large, though robust recoverable capacity will be a small fraction of this theoretical bound.

\subsection{Global Embedding and the Self-Carrying Film}

The most interesting application of hyperpleonastic steganography is global embedding, in which the payload is not a small secret file but structured information about the film itself. Several forms of global embedding are possible.

\textit{Soundtrack embedding.} For a normalized audio signal $a(t)$, the quantized audio amplitude can be mapped to an integer index $i = Q(a_t)$, which selects a realization $r_i \in \mathcal{A}(x)$ at each decision point $x$ that is temporally associated with $a(t)$. The film's material and procedural choices then carry its soundtrack throughout every frame.

\textit{Scene description embedding.} A compressed representation of the world state $\mathcal{W}_t$ (a serialized scene graph excerpt, a procedural seed, or a compressed semantic annotation) is distributed as payload across the admissible differentials of the rendering pipeline. A reader with the decoding scheme can reconstruct the scene description from the rendered video without access to the original world state files.

\textit{Reconstruction information embedding.} The film carries, in its own admissible differentials, sufficient information to reconstruct an alternative representation of itself. This is the strongest form: the film is simultaneously a presentation and a carrier for its own alternative reconstruction.

\subsection{Formal Unification}

The hyperpleonastic steganography extension connects to the broader framework through the following unification:

\[
\text{underdetermination} = \text{creative freedom} = \text{compression opportunity} = \text{steganographic capacity},
\]

subject to the admissibility boundary in all cases. The rendering system that searches for admissible equivalences is simultaneously: identifying degrees of freedom for creative stylistic choices, identifying dimensions that can be aggressively quantized without perceptible loss, and identifying channels through which payload information can be embedded without distorting the visible output.

The admissibility condition is the shared constraint. In all three roles, the selection must remain within $\mathcal{A}(x)$; no realization may be chosen that violates the visible semantic and geometric content of the scene merely to accommodate a creative choice, a compression saving, or a steganographic bit. The penalty for violating admissibility is the same in all three cases: a structural glitch that is detectable as a violation of a recoverable structural invariant.

\subsection{Robustness Note}

A full treatment of robustness, which is the central practical challenge in steganography, is beyond the scope of this appendix. We note only that robustness under re-rendering is qualitatively different from robustness under image processing: a re-rendered film from the same world state $\mathcal{W}_t$ with the same payload will produce the same realization selections, because the encoding scheme is deterministic given $\mathcal{W}_t$ and the payload. Robustness under transcoding by a codec that lacks access to $\mathcal{W}_t$ is a harder problem, since the codec may alter the admissible differentials in ways that corrupt the selection. This is an empirical question to be characterized in the experimental program.

%% ============================================================
\begin{thebibliography}{99}

\bibitem{berger1971}
T.~Berger, \textit{Rate Distortion Theory: A Mathematical Basis for Data Compression}. Prentice-Hall, 1971.

\bibitem{blender2024}
Blender Foundation, \textit{Blender}, version 4.x. \url{https://www.blender.org}, 2024.

\bibitem{coverandthomas2006}
T.~M. Cover and J.~A. Thomas, \textit{Elements of Information Theory}, 2nd ed. Wiley-Interscience, 2006.

\bibitem{davison2007monoslam}
A.~J. Davison, I.~D. Reid, N.~D. Molton, and O.~Stasse, ``MonoSLAM: Real-time single camera SLAM,'' \textit{IEEE Transactions on Pattern Analysis and Machine Intelligence}, vol.~29, no.~6, pp.~1052--1067, 2007.

\bibitem{farneback2003}
G.~Farneb{\"a}ck, ``Two-frame motion estimation based on polynomial expansion,'' in \textit{Scandinavian Conference on Image Analysis}, pp.~363--370, 2003.

\bibitem{ffmpeg2024}
FFmpeg Developers, \textit{FFmpeg}. \url{https://ffmpeg.org}, 2024.

\bibitem{fridrich2009steganography}
J.~Fridrich, \textit{Steganography in Digital Media: Principles, Algorithms, and Applications}. Cambridge University Press, 2009.

\bibitem{hartley2003multiple}
R.~Hartley and A.~Zisserman, \textit{Multiple View Geometry in Computer Vision}, 2nd ed. Cambridge University Press, 2003.

\bibitem{horn1981optical}
B.~K.~P. Horn and B.~G. Schunck, ``Determining optical flow,'' \textit{Artificial Intelligence}, vol.~17, no.~1--3, pp.~185--203, 1981.

\bibitem{kerbl2023gaussians}
B.~Kerbl, G.~Kopanas, T.~Leimk{\"u}hler, and G.~Drettakis, ``3D Gaussian splatting for real-time radiance field rendering,'' \textit{ACM Transactions on Graphics}, vol.~42, no.~4, 2023.

\bibitem{lakshminarayanan2017deep}
B.~Lakshminarayanan, A.~Pritzel, and C.~Blundell, ``Simple and scalable predictive uncertainty estimation using deep ensembles,'' in \textit{Advances in Neural Information Processing Systems}, 2017.

\bibitem{lucas1981optical}
B.~D. Lucas and T.~Kanade, ``An iterative image registration technique with an application to stereo vision,'' in \textit{Proceedings of IJCAI}, pp.~674--679, 1981.

\bibitem{mildenhall2020nerf}
B.~Mildenhall, P.~P. Srinivasan, M.~Tancik, J.~T. Barron, R.~Ramamoorthi, and R.~Ng, ``NeRF: Representing scenes as neural radiance fields for view synthesis,'' in \textit{European Conference on Computer Vision}, 2020.

\bibitem{pharr2023pbr}
M.~Pharr, W.~Jakob, and G.~Humphreys, \textit{Physically Based Rendering: From Theory to Implementation}, 4th ed. MIT Press, 2023.

\bibitem{ranftl2020}
R.~Ranftl, K.~Lasinger, D.~Hafner, K.~Schindler, and V.~Koltun, ``Towards robust monocular depth estimation: Mixing datasets for zero-shot cross-dataset transfer,'' \textit{IEEE Transactions on Pattern Analysis and Machine Intelligence}, vol.~44, no.~3, pp.~1623--1637, 2022.

\bibitem{shannon1948}
C.~E. Shannon, ``A mathematical theory of communication,'' \textit{Bell System Technical Journal}, vol.~27, pp.~379--423 and 623--656, 1948.

\bibitem{sullivan2012}
G.~J. Sullivan, J.-R. Ohm, W.-J. Han, and T.~Wiegand, ``Overview of the High Efficiency Video Coding (HEVC) standard,'' \textit{IEEE Transactions on Circuits and Systems for Video Technology}, vol.~22, no.~12, pp.~1649--1668, 2012.

\bibitem{wang2004ssim}
Z.~Wang, A.~C. Bovik, H.~R. Sheikh, and E.~P. Simoncelli, ``Image quality assessment: From error visibility to structural similarity,'' \textit{IEEE Transactions on Image Processing}, vol.~13, no.~4, pp.~600--612, 2004.

\bibitem{wiegand2003}
T.~Wiegand, G.~J. Sullivan, G.~Bjontegaard, and A.~Luthra, ``Overview of the H.264/AVC video coding standard,'' \textit{IEEE Transactions on Circuits and Systems for Video Technology}, vol.~13, no.~7, pp.~560--576, 2003.

\end{thebibliography}

\end{document}
